% =====================================================================
%  Stage 1, раздел 4.
%  Формализация требований к системе и гипотезы исследования.
% =====================================================================

\section{Требования к системе и формализация гипотезы исследования}

\subsection{Концептуальная модель системы}

В обобщённой постановке разрабатываемая система состоит из трёх подсистем,
взаимодействие которых определено протоколом обмена событиями.

\begin{enumerate}
    \item \textbf{Подсистема предобработки и формирования окон.}
    Принимает поток сетевых соединений в формате IPFIX-совместимых
    flow-записей, выполняет валидацию схемы, очистку, нормализацию и
    формирование скользящих окон фиксированной длины.
    \item \textbf{Подсистема детекции (нейросетевая модель).}
    Принимает окна, формирует скоринг \(s(x)\) и --- после применения
    откалиброванного порога \(\tau\) --- бинарное решение об аномалии,
    сопровождаемое метаданными для приоритизации алертов.
    \item \textbf{Подсистема активного тестирования (ИИ-злоумышленник).}
    Воспроизводит контролируемый поток событий с переключением режимов
    \emph{норма / атака / дрейф} и подаёт его в подсистему предобработки.
\end{enumerate}

В составе SOC рассматриваемая система занимает уровень NTA-аналитики и
взаимодействует с источниками сетевой телеметрии и подсистемой алертирования
(SIEM/SOAR), что согласуется с \secref{subsec:soc-integration} раздела~1.

\subsection{Функциональные требования (\textbf{FR})}

\begin{description}[font=\normalfont\bfseries, leftmargin=2.5em]
    \item[FR-1.] Система должна принимать на вход flow-записи, совместимые
    со схемой UNSW-NB15 (49 признаков) и подмножеством, общим с CICIDS2017
    (для cross-evaluation), с явной валидацией типов и диапазонов.
    \item[FR-2.] Система должна выполнять предобработку признаков:
    очистку пропусков, log-преобразование признаков с положительной
    асимметрией, robust-scaling, кодирование категориальных признаков,
    формирование скользящих окон.
    \item[FR-3.] Система должна обеспечивать обучение, валидацию и сохранение
    весов нейросетевой модели CNN--LSTM (основной) и моделей-альтернатив
    (1D-CNN, LSTM, BiLSTM, GRU, TCN, Transformer-encoder, Autoencoder, VAE,
    MLP) и базовых моделей (Logistic Regression, Random Forest, XGBoost,
    Isolation Forest, One-Class SVM).
    \item[FR-4.] Система должна формировать скоринг \(s(x)\) и применять
    откалиброванный порог \(\tau\), полученный по целевому уровню FP-rate
    или FP-per-time.
    \item[FR-5.] Система должна формировать алерты с метаданными:
    идентификатор окна, метка времени, скоринг, предсказанный класс,
    уровень приоритета, признаковый профиль для аналитика.
    \item[FR-6.] Система должна предоставлять REST-интерфейс на базе FastAPI
    для подачи окна (или последовательности окон) и получения скоринга,
    решения и метаданных.
    \item[FR-7.] Система должна предоставлять минимальный пользовательский
    интерфейс на базе Streamlit для отображения потока алертов, метрик
    качества и состояния drift-мониторинга.
    \item[FR-8.] Система должна включать ИИ-злоумышленника на основе
    \emph{шаблонов + конечного автомата + drift-инжектора} и обеспечивать
    конфигурируемые сценарии тестирования.
    \item[FR-9.] Система должна реализовывать drift-мониторинг по индексам
    PSI, KL и MMD относительно эталонного окна нормы.
    \item[FR-10.] Все эксперименты должны фиксировать seed
    (\(\geq 5\) независимых запусков) и записывать метрики и
    конфигурации в журналы (CSV/JSON) для воспроизведения.
\end{description}

\subsection{Нефункциональные требования (\textbf{NFR})}

\begin{description}[font=\normalfont\bfseries, leftmargin=2.5em]
    \item[NFR-1.] \emph{Качество детекции на UNSW-NB15.}
    F1 \(\geq 0{,}90\) и PR-AUC \(\geq 0{,}95\) на официальном
    тестовом разделении UNSW-NB15 при бинарной постановке.
    \item[NFR-2.] \emph{Уровень ложных срабатываний.} FPR \(\leq 1\%\)
    при Recall \(\geq 0{,}85\) после калибровки порога; контроль
    FP-per-time для оценки эксплуатационной нагрузки.
    \item[NFR-3.] \emph{Эксплуатационная задержка.} Inference latency
    \(\leq 50\) мс на одно окно при batch-size = 1 на CPU-конфигурации
    среднего класса.
    \item[NFR-4.] \emph{Пропускная способность.} Throughput
    \(\geq 1000\) flow-записей в секунду (EPS) при батчировании
    окон на CPU-конфигурации среднего класса.
    \item[NFR-5.] \emph{Time-to-detect.} TTD \(\leq 1\) окна с момента
    появления атаки в потоке (для шаблонов с явными flow-проявлениями
    --- DDoS, Port Scan, Brute Force).
    \item[NFR-6.] \emph{Устойчивость к дрейфу.} Снижение F1 при включении
    drift-инжектора с интенсивностью PSI~\(\leq 0{,}25\) не должно
    превышать абсолютной величины \(\Delta\,\text{F1} = 0{,}10\) до
    срабатывания механизма адаптации/повторной калибровки.
    \item[NFR-7.] \emph{Перенос на CICIDS2017.} В режиме \emph{train on
    UNSW, test on CICIDS2017} (без дообучения) \(\Delta\,\text{F1}\)
    относительно качества на UNSW \(\leq 0{,}20\); с дообучением на 10\%
    выборки CICIDS2017 \(\Delta\,\text{F1} \leq 0{,}10\).
    \item[NFR-8.] \emph{Воспроизводимость.} Полный pipeline (от raw-CSV
    до итоговых метрик и графиков отчёта) запускается одной командой
    Makefile при фиксированных seed.
    \item[NFR-9.] \emph{Стабильность.} Дисперсия F1 по seed
    (\(n \geq 5\)) удовлетворяет \(\sigma_{F1} \leq 0{,}01\).
\end{description}

\subsection{Метрики и протокол оценки}

В соответствии с \secref{sec:metrics} раздела~1, для оценки разработанной
методики применяются три группы метрик.

\begin{enumerate}
    \item \textbf{Классификационные:} Precision, Recall, F1, MCC, FPR.
    \item \textbf{Интегральные:} PR-AUC (приоритетно), ROC-AUC.
    \item \textbf{Эксплуатационные:} time-to-detect, latency, throughput,
    FP-per-time.
\end{enumerate}

Дополнительно фиксируются:
\begin{itemize}
    \item Калибровка предсказаний (ECE);
    \item Метрики дрейфа (PSI, MMD относительно эталонного окна);
    \item Стабильность по seed (\(\sigma_{F1}\), \(\sigma_{\text{PR-AUC}}\)).
\end{itemize}

Протокол сравнения моделей: одинаковое разделение, одинаковый препроцессинг,
одинаковый порог калибровки (по целевому FP-rate). Статистическая значимость
различий между моделями оценивается парным критерием Уилкоксона по
результатам \(n \geq 5\) запусков.

\subsection{Допущения и ограничения исследования}

\begin{enumerate}
    \item \textbf{Допущение об отсутствии deep packet inspection.} Анализ
    выполняется на flow-уровне; полезная нагрузка пакетов не учитывается.
    Такое допущение согласуется с реальной практикой применения NIDS в
    шифрованном трафике и со стандартом IPFIX~\cite{rfc7011}.
    \item \textbf{Допущение об автономности активного тестирования.}
    ИИ-злоумышленник работает в рамках программного стенда и не порождает
    реальный сетевой трафик за пределами стенда.
    \item \textbf{Ограничение по числу классов.} Мультиклассовая постановка
    рассматривается дополнительно (см. главу~2), однако основная гипотеза
    проверяется в бинарной постановке \emph{normal/attack}.
    \item \textbf{Ограничение по типам дрейфа.} Моделируются три типа
    дрейфа (covariate, prior, concept). Recurring drift и более сложные
    нестационарные режимы отнесены к направлениям дальнейшего развития.
    \item \textbf{Ограничение по DL-фреймворку.} Модели реализуются на
    PyTorch, классические baseline --- на scikit-learn и XGBoost. Сравнение
    с другими фреймворками вне области исследования.
\end{enumerate}

\subsection{Формализация гипотезы исследования}

\subsubsection{Основная гипотеза}

\begin{statement}[Основная гипотеза, $H_0^{\text{main}}$]
Существует нейросетевая модель \(M^*\) с архитектурой \emph{CNN--LSTM},
обученная на flow-представлении датасета UNSW-NB15 в скользящих окнах,
такая, что на официальном тестовом разделении датасета:
\begin{enumerate}
    \item F1 \(\geq 0{,}90\) и PR-AUC \(\geq 0{,}95\),
    \item FPR \(\leq 1\%\) при Recall \(\geq 0{,}85\),
\end{enumerate}
и эти показатели статистически значимо превосходят (\(p < 0{,}05\) по парному
критерию Уилкоксона) показатели ансамбля моделей-baseline (Logistic
Regression, Random Forest, XGBoost, Isolation Forest, One-Class SVM, MLP)
при равных условиях обучения и оценки.
\end{statement}

\subsubsection{Дополнительная гипотеза о роли активного тестирования}

\begin{statement}[$H_0^{\text{aux}}$]
Применение ИИ-злоумышленника в форме \emph{шаблоны + конечный автомат +
drift-инжектор} к разработанному прототипу позволяет:
\begin{enumerate}
    \item обнаруживать аномалии в потоковом сценарии с TTD \(\leq 1\) окна
    для шаблонных атак с явными flow-проявлениями (DDoS, Port Scan, Brute
    Force);
    \item фиксировать снижение качества детекции при контролируемом
    дрейфе (PSI~\(\leq 0{,}25\)) в пределах \(\Delta\,\text{F1} \leq 0{,}10\);
    \item систематически идентифицировать классы со сниженной полнотой
    обнаружения и связывать ошибки FN с конкретными состояниями FSM, что
    обеспечивает дифференциальный анализ устойчивости детектора, недоступный
    при использовании только статического датасета.
\end{enumerate}
\end{statement}

\subsubsection{Гипотеза переносимости}

\begin{statement}[$H_0^{\text{transfer}}$]
Модель \(M^*\), обученная на UNSW-NB15, демонстрирует на CICIDS2017
\(\Delta\,\text{F1} \leq 0{,}20\) без дообучения и
\(\Delta\,\text{F1} \leq 0{,}10\) с дообучением на 10\% выборки CICIDS2017
относительно качества на UNSW-NB15.
\end{statement}

\subsection{Критерии подтверждения и опровержения гипотезы}

Гипотеза \(H_0^{\text{main}}\) считается подтверждённой, если на
\(n \geq 5\) запусках с независимыми seed выполнены условия NFR-1, NFR-2 и
получено статистически значимое превосходство над ансамблем baseline
(\(p < 0{,}05\)).

Гипотеза \(H_0^{\text{aux}}\) подтверждается при выполнении NFR-5 для
шаблонных атак, NFR-6 для drift-сценариев и при наличии в результатах
систематического разделения ошибок FN по состояниям FSM.

Гипотеза \(H_0^{\text{transfer}}\) подтверждается при выполнении NFR-7.

В случае невыполнения какого-либо из критериев гипотеза подвергается
уточнению по итогам анализа: возможны изменения архитектуры (переход к
BiLSTM~+~Attention, Transformer-encoder), параметров препроцессинга,
порога калибровки или политики FSM. Каждое такое изменение фиксируется в
журнале экспериментов.

\subsection{Промежуточные выводы по разделу}

\begin{enumerate}
    \item Сформирована концептуальная трёхуровневая модель системы
    (предобработка, детекция, активное тестирование) с явным определением
    интерфейсов и места в структуре SOC.
    \item Определены десять функциональных и девять нефункциональных
    требований с количественными порогами по качеству детекции, задержкам,
    устойчивости к дрейфу и переносимости. Все пороги связаны с метриками
    \secref{sec:metrics} и обоснованы с учётом эксплуатационной природы
    NIDS.
    \item Сформулирована основная гипотеза исследования \(H_0^{\text{main}}\)
    о превосходстве CNN--LSTM на UNSW-NB15 над ансамблем baseline, а также
    дополнительные гипотезы о роли ИИ-злоумышленника
    (\(H_0^{\text{aux}}\)) и о переносимости методики на CICIDS2017
    (\(H_0^{\text{transfer}}\)) с явными критериями подтверждения и
    опровержения.
    \item Определены допущения и ограничения исследования (flow-уровень
    анализа, программно-стендовое тестирование, бинарная постановка как
    основная), что задаёт область применимости разработанной методики и
    границы интерпретации экспериментальных результатов.
\end{enumerate}
